% Draft status: V0
% Evidence status: checked where possible; TODOs mark missing evidence.
\section{Experiments}

The experiments compare ImageNet-pretrained evaluation with dataset-specific fine-tuning. In the pretrained state, each model is evaluated on the processed GeoDE or Dollar Street test split before dataset-specific training. In the fine-tuned state, the model is initialized from pretrained weights, fine-tuned on the dataset training split, and evaluated on the held-out test split. Results are reported both overall and by region.

The processed datasets use train, validation, and test directories. The local class mappings contain 40 classes for GeoDE and 64 classes for the processed Dollar Street subset. The Dollar Street subset is therefore not the full official 289-category Dollar Street label space; it is an ImageNet-compatible experimental subset. The pretrained evaluation files record that, when the number of target classes differs from 1000, the evaluation uses pretrained weights where shapes match plus a newly initialized dataset classifier head. This detail is important when interpreting the very low pretrained accuracies.

The fine-tuning runs use a common protocol recorded in the local configuration files: input image size 224, batch size 32, Adam optimization, learning rate \(3\times10^{-5}\), cosine scheduling, label smoothing of 0.1, automatic mixed precision, and 10 epochs. Training uses the dataset training split, model selection uses validation performance, and final metrics are computed on the test split. The evaluation script computes object accuracy, region-conditioned object accuracy, best region, worst region, and the best-minus-worst geographic-disparity gap.

The model initialization follows the intended taxonomy in Table~\ref{tab:revision-model-taxonomy}. The external training script maps the experimental model keys to pretrained model identifiers for ResNet50, ConvNeXtV2, DeiT3, SwinV2, MaxViT, DINOv2, MAE, CLIP, and SigLIP. TODO: before final submission, audit the Dollar Street SwinV2 configuration mismatch noted in the Models section and verify that all result rows correspond to the intended pretrained model identifiers.

All metrics used in the paper are normalized into \texttt{derived/processed/all\_metrics.csv}. The table contains 36 rows: two datasets, nine models, and two training states. The validation scripts check that all expected rows are present, that fine-tuning state values are valid, and that geographic disparity equals best-region accuracy minus worst-region accuracy. No model training is performed inside the paper workspace.
